% TOP_PROBLEM: {"problemId": "problem.exact-real-grothendieck-constant", "canonicalTitle": "Exact Real Grothendieck Constant", "releaseRank": 310, "primaryDomain": "analysis_pde", "primaryDomainLabel": "Analysis & PDE", "displayStatus": "open", "releaseStatus": "open"}
% !TeX program = lualatex
% ATTEMPT_STATUS: unresolved
% Model: GPT-6 Astra Ultra; continuation dated 27 September 2026.
\documentclass[11pt]{article}
\usepackage[margin=25mm]{geometry}
\usepackage{fontspec}
\setmainfont{Latin Modern Roman}
\usepackage{amsmath,amssymb,mathtools}
\usepackage{unicode-math}
\setmathfont{Latin Modern Math}
\usepackage{xurl,hyperref}
\hypersetup{colorlinks=true,urlcolor=blue}
\setcounter{secnumdepth}{0}
\setlength{\parindent}{0pt}
\setlength{\parskip}{5pt}
\setlength{\emergencystretch}{3em}
\begin{document}
\section*{\texorpdfstring{Exact Real Grothendieck Constant}{Exact Real Grothendieck Constant}}\label{310}
\subsection{Definitions and mathematical statement}
For a real $m\times n$ matrix $A=(a_{ij})$, put
\[
 B(A)=\max_{{\varepsilon}_i,\delta_j\in\{-1,1\}}
 \left|\sum_{i,j}a_{ij}{\varepsilon}_i\delta_j\right|.
\]
The real Grothendieck constant is the least $K$ such that, for every $m,n,A$ and all unit vectors $u_i,v_j$ in any real Hilbert space,
\[
 \left|\sum_{i,j}a_{ij}\langle u_i,v_j\rangle\right|\le K B(A).
\]
Determine this optimal universal constant $K_G^{{\mathbb{R}}}$ exactly. The dimensions and the matrix are arbitrary, and the comparison is with independent signs on rows and columns. Neither a fixed-dimensional variant nor the complex Grothendieck constant is the selected quantity. An improved upper or lower estimate alone is not an exact determination.
\par\smallskip\textit{Formulation and concept references:} \href{https://arxiv.org/abs/1103.6161}{[S1]}.

\subsection{Short English statement}
What is the exact largest advantage that real unit-vector inner products can have over signs in a bilinear form?
\subsection{Research attempt}
\paragraph{The universal optimization and its status.}
The exact real constant is \textbf{UNRESOLVED} by this attempt.
The primary source [S1] proves a strict improvement over Krivine's
classical upper bound. Accordingly, the elementary derivation of
that bound below is a valid certificate, not a proposed exact answer
and not the best known upper estimate. The useful outputs are an
exact small-matrix lower bound, a complete Gaussian rounding proof,
and checks on two tempting but invalid amplification arguments.

Write $V(A)$ for the supremum of the absolute vector expression
in the question. Simultaneously changing all row vectors to their
negatives changes its sign, so absolute values can be replaced by
maximization of the signed expression. The same applies to $B(A)$.
If $A\ne0$, then $B(A)>0$: averaging the square of the signed sum
over independent uniform signs gives $\sum_{ij}a_{ij}^2>0$.
Consequently
\[
 K_G^{\mathbb R}=\sup_{m,n,A\ne0}\frac{V(A)}{B(A)}.
 \tag{310.1}
\]
This is a supremum over arbitrary finite sizes, not the optimization
of one prescribed matrix.

\paragraph{Exact evaluation of the sign side.}
For fixed column signs $\delta$, choose each row sign to agree
with the sign of $\sum_j a_{ij}\delta_j$, choosing either sign
when that sum is zero. This proves
\[
 B(A)=\max_{\delta\in\{-1,1\}^n}
              \sum_i\left|\sum_j a_{ij}\delta_j\right|.
 \tag{310.2}
\]
It gives an exact finite algorithm for rational matrices, albeit
an exponential one. It also shows that imposing the same sign on
row and column copies would change the problem: their choices are
independent, even for a square symmetric matrix.

For fixed $m,n$, the vector side needs no Hilbert space of dimension
larger than $m+n$. Only the span of the finitely many vectors
matters. Equivalently introduce their Gram matrix $G$; the feasible
set is
\[
 G\succeq0,\qquad G_{aa}=1\quad(1\leq a\leq m+n).
 \tag{310.3}
\]
The objective is the linear functional
$\sum_{ij}a_{ij}G_{i,m+j}$. This set is closed and bounded because
positive semidefiniteness gives $|G_{ab}|\leq1$, so the maximum is
attained. Conversely every such matrix has a real Gram factorization.
This finite-dimensional compactness does not bound the sizes needed
in the outer supremum (310.1).

\paragraph{A rigorous lower bound from a two-by-two matrix.}
Take
\[
 H=\begin{pmatrix}1&1\\1&-1\end{pmatrix}.
\]
Equation (310.2) gives $B(H)=2$: of $\delta_1+\delta_2$ and
$\delta_1-\delta_2$, exactly one has absolute value two.
For arbitrary unit $v_1,v_2$, optimizing the row vectors gives
\[
 V(H)=\max\bigl(\|v_1+v_2\|+\|v_1-v_2\|\bigr)=2\sqrt2.
 \tag{310.4}
\]
Indeed the sum of the squares of these two norms is four, so
Cauchy--Schwarz bounds their sum by $2\sqrt2$. Equality holds for
orthogonal $v_1,v_2$, with row vectors along their sum and difference.
Thus $K_G^{\mathbb R}\geq\sqrt2$. This is an exact example in a
two-dimensional real space, not an assertion that the universal
constant equals its dimension-two variant.

\paragraph{Gaussian signs and the arcsine identity.}
Let $x,y$ be unit vectors with $\langle x,y\rangle=t$, and let
$Z$ be a standard Gaussian in their finite-dimensional span. Then
\[
 \mathbb E\bigl[\operatorname{sgn}\langle Z,x\rangle
                 \operatorname{sgn}\langle Z,y\rangle\bigr]
       =\frac2\pi\arcsin t.
 \tag{310.5}
\]
For a proof, write $t=\cos\theta$ with $0\leq\theta\leq\pi$.
The angular part of a rotationally invariant two-dimensional
Gaussian is uniform. The two scalar products have opposite signs
on sectors of total angular measure $2\theta$, hence with probability
$\theta/\pi$. Their product expectation is $1-2\theta/\pi$,
which is the right side. The collinear endpoints follow by direct
inspection or continuity; zero scalar products have probability zero.

Applying (310.5) directly rounds correlations to arcsines, not to
a fixed multiple of themselves. One cannot multiply a termwise
inequality by arbitrary signed coefficients $a_{ij}$ and sum it.
The next construction removes that obstruction exactly.

\paragraph{Tensor embeddings linearize the rounding.}
Put
\[
 c=\operatorname{arsinh}(1)=\log(1+\sqrt2),\qquad
 b_r=\frac{c^{2r+1}}{(2r+1)!}\quad(r\geq0).
\]
For a real unit vector $u$, define vectors in the Hilbert direct sum
of odd tensor powers by
\[
 U(u)=\bigoplus_{r\geq0}\sqrt{b_r}\,u^{\otimes(2r+1)},
 \qquad
 W(v)=\bigoplus_{r\geq0}(-1)^r\sqrt{b_r}\,
                                      v^{\otimes(2r+1)}.
 \tag{310.6}
\]
The series converge because $\sum_rb_r=\sinh c=1$.
Both vectors have norm one, and the tensor inner-product rule gives
\[
 \langle U(u),W(v)\rangle
 =\sum_{r\geq0}(-1)^rb_r\langle u,v\rangle^{2r+1}
 =\sin(c\langle u,v\rangle).
 \tag{310.7}
\]
Although the ambient sum is infinite-dimensional, the finitely many
$U(u_i),W(v_j)$ span a finite-dimensional space. A standard Gaussian
on that span therefore suffices; no Gaussian vector with identity
covariance on an infinite-dimensional Hilbert space is assumed.

Use one Gaussian to round all these transformed vectors to signs.
Since $c<\pi/2$ and $|\langle u,v\rangle|\leq1$,
$\arcsin(\sin(ct))=ct$ on the entire required interval. Thus
\[
 \mathbb E\sum_{ij}a_{ij}\varepsilon_i\delta_j
       =\frac{2c}{\pi}\sum_{ij}a_{ij}\langle u_i,v_j\rangle.
\]
Every realized sign sum has absolute value at most $B(A)$.
Taking absolute expectations proves the uniform inequality
\[
 V(A)\leq\frac{\pi}{2\log(1+\sqrt2)}B(A).
 \tag{310.8}
\]
All signs of the coefficients have been retained. Together with
(310.4), this proves a nontrivial finite interval for the universal
constant by explicit constructions.

\paragraph{What optimizing this construction actually establishes.}
For the same sine tensor prescription with another positive number
$c'$, the two squared norms are $\sinh c'$. To keep them at most
one before adding orthogonal padding, one needs $c'\leq c$. Hence
$c$ is optimal for this particular coefficient prescription. This
does not optimize every possible embedding or correlated rounding
scheme. The source's strict improvement is precisely a warning
against promoting an internally optimized method into the optimal
universal constant.

Similarly, (310.4) saturates two elementary inequalities for $H$.
That says nothing about a matrix with many more rows and columns.
There is no principle in (310.1) that an extremizer must occur in
the smallest nontrivial size. To prove an exact value one needs a
universal upper bound and matching examples approaching it, with
the same real scalar and independent-sign conventions.

\paragraph{Direct sums and tensor products fail to amplify naively.}
For block diagonal matrices $A\oplus B$, independent row and column
choices on the blocks give
\[
 B(A\oplus B)=B(A)+B(B),\qquad
 V(A\oplus B)=V(A)+V(B).
\]
The upper inequalities follow by bounding each block separately;
the lower ones use optimal choices for each block, changing signs
so their contributions are positive. Therefore the ratio for a
direct sum is a weighted average of the two ratios, not their sum
or product.

Tensor products permit signs indexed by pairs; these need not split
into products of signs on individual factors. This invalidates the
claim $B(A\otimes B)=B(A)B(B)$. A complete countercheck is
$H\otimes H$, a four-by-four Hadamard matrix $Q$ with
$QQ^{\mathsf T}=4I$. For every sign vector $\delta$,
\[
 \|Q\delta\|_1\leq2\|Q\delta\|_2=8.
\]
For $\delta=(1,1,1,-1)$, all four entries of $Q\delta$ have
absolute value two, so $B(Q)=8$, whereas $B(H)^2=4$.
In fact $V(Q)=8$ as well. For arbitrary unit column vectors,
Cauchy--Schwarz and orthogonality of the columns of $Q$ give
\[
 \sum_{i=1}^4\Big\|\sum_{j=1}^4Q_{ij}v_j\Big\|
 \leq2\left(\sum_i\Big\|\sum_jQ_{ij}v_j\Big\|^2\right)^{1/2}
 =8.
\]
The sign example attains it. Tensoring this lower-bound witness
therefore reduces its ratio to one, rather than squaring $\sqrt2$.

\paragraph{Verification and the remaining gap.}
The saved finite checks enumerate signs exactly for $H$ and
$H\otimes H$, verify their integer orthogonality identities, and
numerically compare truncated tensor series and the arcsine formula
with the stated elementary functions. No sampled Gaussian average
is used in place of the exact angular proof of (310.5).

The proof of (310.8) is complete but its constant is known from
[S1] to be nonoptimal. The first unresolved step is a sharp rounding
or dual inequality together with matching unrestricted matrix
examples. Finite semidefinite programs, direct sums and tensoring
the displayed witness do not supply either half of that equality.
The exact value in (310.1) remains \textbf{UNRESOLVED} here.
\subsection{Sources}
\begin{itemize}
\item[S1]1103.6161v3, The Grothendieck constant is strictly smaller than Krivine's bound.
\url{https://arxiv.org/abs/1103.6161}.
\textit{Formulation source.}
\end{itemize}
\end{document}
